\chapter{Modelos Comentados dos Elementos}

Este capítulo apresenta modelos esquemáticos dos principais elementos do TCC. Os modelos não substituem o padrão gráfico oficial, mas ajudam o estudante a compreender a função de cada parte e a conferir se os dados essenciais foram preenchidos.

\section{Capa}

\modeloinstitucional{MODELO ESQUEMÁTICO DE CAPA}{
\centering
UNIVERSIDADE DO DISTRITO FEDERAL PROFESSOR JORGE AMAURY MAIA NUNES -- UnDF\par
\vspace{2em}
NOME DO AUTOR\par
\vspace{4em}
TÍTULO DO TRABALHO\par
Subtítulo, se houver\par
\vfill
BRASÍLIA\par
ANO\par
}

A capa deve identificar o trabalho de modo direto. O título precisa ser fiel ao conteúdo e não deve prometer mais do que a pesquisa entrega. O subtítulo, quando usado, deve complementar o título, não repeti-lo.

Erros comuns: título genérico, nome do autor diferente da folha de rosto, ano desatualizado, excesso de informação na capa e uso de siglas não explicadas.

\section{Folha de rosto}

\modeloinstitucional{MODELO ESQUEMÁTICO DE FOLHA DE ROSTO}{
\centering
NOME DO AUTOR\par
\vspace{3em}
TÍTULO DO TRABALHO\par
Subtítulo, se houver\par
\vspace{3em}
\raggedleft
\begin{minipage}{0.58\linewidth}
Trabalho de Conclusão de Curso apresentado ao curso de [CURSO] da Universidade do Distrito Federal Professor Jorge Amaury Maia Nunes -- UnDF, como requisito parcial para obtenção do título de [GRAU].\par
\vspace{0.5em}
Orientador: Nome do Orientador.
\end{minipage}\par
\vfill
\centering
BRASÍLIA\par
ANO\par
}

A folha de rosto repete dados centrais e acrescenta a natureza do trabalho. É nela que o leitor identifica o curso, o grau pretendido e a orientação.

\section{Folha de aprovação}

\modeloinstitucional{MODELO ESQUEMÁTICO DE FOLHA DE APROVAÇÃO}{
\centering
NOME DO AUTOR\par
\vspace{2em}
TÍTULO DO TRABALHO\par
\vspace{2em}
\raggedright
Trabalho aprovado como requisito parcial para obtenção do título de [GRAU] no curso de [CURSO] da Universidade do Distrito Federal Professor Jorge Amaury Maia Nunes -- UnDF.\par
\vspace{1.2em}
Brasília, DF, dia de mês de ano.\par
\vspace{2em}
\centering
\rule{0.68\linewidth}{0.4pt}\par
Prof. Dr. Nome do Orientador\par
Instituição\par
\vspace{1.5em}
\rule{0.68\linewidth}{0.4pt}\par
Profa. Dra. Nome do Membro\par
Instituição\par
}

A folha de aprovação deve manter a banca em formato uniforme. Quando o padrão institucional adota titulação abreviada antes do nome, ela não precisa ser repetida em linha separada.

\section{Errata}

\modeloinstitucional{MODELO ESQUEMÁTICO DE ERRATA}{
\centering
ERRATA\par
\vspace{1em}
\raggedright
SOBRENOME, Nome. Título do trabalho. Ano. Trabalho de Conclusão de Curso (Graduação em Curso) -- Universidade do Distrito Federal Professor Jorge Amaury Maia Nunes -- UnDF, Brasília, ano.\par
\vspace{1em}
\centering
\begin{tabular}{cccc}
\toprule
Folha & Linha & Onde se lê & Leia-se \\
\midrule
16 & 10 & auto-clavado & autoclavado \\
\bottomrule
\end{tabular}
}

A errata deve ser usada apenas para correções pontuais. Se o problema exige reescrever seções inteiras, o correto é revisar a versão final antes da entrega.

\section{Resumo}

\modeloinstitucional{MODELO ESQUEMÁTICO DE RESUMO}{
\centering
RESUMO\par
\vspace{1em}
\raggedright
Texto conciso apresentando objetivo, método, resultados e conclusões do trabalho. O resumo deve permitir que o leitor compreenda a essência da pesquisa sem consultar o texto integral.\par
\vspace{1em}
\textbf{Palavras-chave:} termo 1; termo 2; termo 3.
}

O resumo deve evitar citações, notas, ilustrações e informações não desenvolvidas no trabalho. As palavras-chave precisam representar conceitos centrais.

\section{Sumário}

\modeloinstitucional{MODELO ESQUEMÁTICO DE SUMÁRIO}{
\raggedright
1 INTRODUÇÃO \dotfill 10\par
2 REFERENCIAL TEÓRICO \dotfill 15\par
3 METODOLOGIA \dotfill 28\par
4 RESULTADOS E DISCUSSÃO \dotfill 35\par
5 CONSIDERAÇÕES FINAIS \dotfill 48\par
REFERÊNCIAS \dotfill 52\par
APÊNDICE A -- Roteiro de entrevista \dotfill 56\par
ANEXO A -- Documento institucional \dotfill 58\par
}

O sumário deve refletir exatamente os títulos e páginas da versão final. Elementos pré-textuais não entram no sumário.

\section{Referências}

\modeloinstitucional{MODELO ESQUEMÁTICO DE REFERÊNCIAS}{
\raggedright
SOBRENOME, Nome. \textbf{Título da obra}. Edição. Local: Editora, ano.\par
\vspace{0.5em}
SOBRENOME, Nome. Título do artigo. \textbf{Título do periódico}, local, volume, número, páginas, ano. Disponível em: endereço. Acesso em: dia mês ano.\par
}

As referências devem ser conferidas uma a uma. A lista final deve ser coerente com as citações do texto.

\section{Apêndice e anexo}

\modeloinstitucional{MODELO ESQUEMÁTICO DE APÊNDICE E ANEXO}{
\centering
APÊNDICE A -- ROTEIRO DE ENTREVISTA\par
\vspace{2em}
\raggedright
Material elaborado pelo autor do trabalho.\par
\vspace{2em}
\centering
ANEXO A -- DOCUMENTO INSTITUCIONAL\par
\vspace{2em}
\raggedright
Material não elaborado pelo autor, mas necessário à compreensão do trabalho.
}

Apêndices e anexos devem ser citados no corpo do texto. Material não citado tende a parecer acréscimo sem função acadêmica.

